\documentclass[10pt,a4paper]{amsart}
\usepackage[margin=19mm]{geometry}
\usepackage{amsmath,amssymb,mathtools}
\usepackage[scr=rsfs,cal=euler]{mathalfa}
\usepackage{xfrac}
\usepackage[unicode,hidelinks,bookmarks=false]{hyperref}
\newcommand{\ie}{i.e.\ }
\newcommand{\cf}{cf.\ }
\newcommand{\resp}{respectively\ }
\newcommand{\Subsection}{\subsection}
\providecommand{\nobreakdash}{\nobreak\hbox{-}}
\setlength{\emergencystretch}{2em}
\multlinegap0pt
\usepackage{bm}
\usepackage{bbm}
\usepackage{dsfont}
\usepackage{xspace}
\usepackage{cancel}
\usepackage{mathtools}
\usepackage{amsthm}
\makeatletter
\@ifundefined{lemma}{\newtheorem{lemma}{Lemma}}{}
\makeatother
\newcommand{\Pof}[1]{\mathbb{P}\left(#1\right)}
\title[Dimension drop of harmonic measure for some finite range ran]{Dimension drop of harmonic measure for some finite range random walks on Fuchsian Schottky groups}
\author{Ernesto Garc\'ia, Pablo Lessa}
\date{Source: arXiv:2301.09714v3 (CC BY 4.0)}
\begin{document}
\maketitle

\noindent\emph{Source and licence.} Dimension drop of harmonic measure for some finite range random walks on Fuchsian Schottky groups. Version arXiv:2301.09714v3 (CC BY 4.0), distributed under the Creative Commons Attribution 4.0 International licence, as recorded in the arXiv licence metadata for this exact version and shown on the abstract page https://arxiv.org/abs/2301.09714v3. Full rights notice: licenses/arxiv-2301.09714-CC-BY-4.0.txt.\par\smallskip

\noindent\textit{Editorial scope.} Counted unit: the Lemma whose source label is \texttt{markovlemma} in the article (source0.tex lines 343--351, source label \texttt{markovlemma}), delivered together with the article's own complete original proof of it (source0.tex lines 352--386, one \texttt{proof} environment of 35 lines). No mathematics is written, repaired or re-derived: statement and proof are the article's own text line for line. Required context retained uncounted: none. Every definition, display and label of those lines is retained unchanged. Omitted from this package and not redistributed: 1--342, 387--566. Nothing inside a retained excerpt is omitted or altered: every retained body is delivered exactly as the article's lines are written, so no omission receipt of any kind accompanies this package. Citations: the retained excerpts carry no citation command at all, so no bibliography entry of the article is transcribed and no reference is redistributed. Reference adaptation: none was needed and none was made; every reference inside the delivered unit resolves inside it, so no unresolved reference or citation is produced. Printed slips kept literal: no printed word, symbol, hypothesis or number of the counted statement or of its proof was altered. Layout and class: the article's own class is \texttt{article} with packages including inputenc, amsmath, amssymb, amsthm, hyperref, cleveref, graphicx, float, bbm, xcolor; the wrapper is the lane's pinned \texttt{amsart} editorial wrapper at 10pt on A4, which restates the theorem-like environments the retained text needs and adds the article's own macro definitions that the retained text needs (\texttt{\textbackslash Pof}), with extra wrapper packages (bm, bbm, dsfont, xspace, cancel, mathtools). The delivered numbering is the unit's own. Assets: no retained span contains a figure environment, a graphics-inclusion command, a PDF-page inclusion, a table or a TikZ picture; no image file is copied into this package, the package declares no asset dependency and no image file is redistributed with it. Licence: version arXiv:2301.09714v3 of this article is distributed under the Creative Commons Attribution 4.0 International licence, as recorded in the arXiv licence metadata for this exact version and shown on the abstract page https://arxiv.org/abs/2301.09714v3. A full rights notice is delivered at licenses/arxiv-2301.09714-CC-BY-4.0.txt. Characters: every retained excerpt is pure ASCII, so the delivered engine, XeLaTeX with its default Latin Modern text font, reports no missing character.
\begin{lemma}\label{markovlemma}
 Let \(r = \lim\limits_{n \to +\infty}\overline{\pi(y_1)\cdots \pi(y_n)}\) and set \(\tau_n = \min\lbrace k: \overline{\pi(y_1)\cdots \pi(y_k)} = r[1:n]\rbrace\) for \(n = 1,2,\ldots\).
 
 Then the sequence \(x_1 = y_{\tau_1},x_2 = y_{\tau_2},\ldots\) is a Markov chain and furthermore
 \begin{equation}\label{markovequation}
 \Pof{x_1 = s_1,\ldots, x_n = s_n} = f(s_1)g(s_1,s_2)\cdots g(s_{n-1},s_n)h(s_n),
 \end{equation}
 for all \(s_1,\ldots, s_n \in S\).
\end{lemma}

\begin{proof}
It suffices to prove equation \ref{markovequation}, since this implies that \(x_1,\ldots, x_n\) is Markov with initial distribution \(p(s) = f(s)h(s)\) and transtion probabilities \(P(s,s') = h(s)^{-1}g(s,s')h(s')\).

For this purpose we let \(w = \pi(s_1)\cdots \pi(s_n)\) and write
\begin{equation} 
\Pof{x_1 = s_1,\ldots, x_n = s_n} = \Pof{F_1 \cap \cdots F_n \cap G} = \Pof{F_n \cap G},
\end{equation}
where \(F_i = \lbrace \tau_{w[1:k]} < +\infty, y_{\tau_{w[1:k]}} = s_k\text{ for all }k \le i\rbrace\) and \(G = \lbrace r\text{ reduced word}: w \prec r\rbrace\).

We observe that conditioned on \(F_n\) the event \(G\) is equivalent to
\[\left\lbrace\pi(s_n) \prec \overline{\pi(y_{\tau_w})\cdots \pi(y_{k + \tau_w})}\text{ for all }k\text{ large enough}\right\rbrace.\]

Therefore, by the strong Markov property we obtain 
\begin{equation}\Pof{F_n \cap G} = \Pof{F_n}\Pof{G | F_n} = \Pof{F_n}h(s_n).
\end{equation}


For each \(i = 1,\ldots,n\) we observe that conditioned on \(F_i\) the event \(F_{i+1}\) is equivalent to there being some finite \(k\) such that \(\overline{\pi(y_{\tau_{w[1:i]}})\cdots \pi(y_{k + \tau_{w[1:i]}})} = \pi(s_i)\pi(s_{i+1})\) and the minimal \(k\) with this property satisfying \(y_{k + \tau_{w[1:i]}} = s_{i+1}\).

Once again by the strong Markov property we obtain 
\begin{equation}
\Pof{F_{i+1}} = \Pof{F_i}\Pof{F_{i+1}|F_i} = \Pof{F_i}g(s_i,s_{i+1}).
\end{equation}


Hence we have 
\begin{align*}
 \Pof{x_1 = s_1,\ldots, x_n = s_n} &= \Pof{F_n \cap G}
 \\ &= \Pof{F_n}\Pof{G|F_n}
 \\ &= \Pof{F_{n-1}}\Pof{F_n|F_{n-1}}h(s_n)
 \\ &= \cdots 
 \\ &= \Pof{F_1}g(s_1,s_2)\cdots g(s_{n-1},s_n)h(s_n),
\end{align*}
the equality \(\Pof{F_1} = f(s_1)\) is direct from the definition of \(f\), which concludes the proof.
\end{proof}

\end{document}
